\scanpage{368}
\mspage{485}

\uline{$\overline{\mathrm{VIII}}$}\quad
\uline{Recapitulation on the equation}
\[
\tag{1}
[\alpha,\rho]=[\beta,\sigma]\varepsilon_1(\nu).
\]
The groups \deleted{$\unsure{\mathcal G_{\rho,\sigma}^{*0}}$}, $S\mathcal G_{\rho,\sigma}^{*}$ and its variants.

The essential result we have arrived at in (\uline{$\overline{\mathrm{III}}$}), is the possibility of parametrizing the solutions of such an equation \inserted{in $\alpha,\beta,\nu$} (under suitable hypotheses on $\rho,\sigma$) by a group $B'_0\times T_\rho\times T_\sigma$, where $B'_0$ is a ``\unsure{pre-Borel|pri-Borel}'' in $G$ (taking, for a suitable basis, the form
\[
B'_0=\left\{\begin{pmatrix}\mu&\lambda\\0&1\end{pmatrix}\;\middle|\;
\begin{array}{l}\mu\in\Gamma \mathbb G_m\\\lambda\in\Gamma \mathbb G_a\end{array}\right.
\]
and $T_\rho,T_\sigma$ are the \unsure{centralizers|stabilizers} of the regular sections $\rho,\sigma$ of $G$, by the morphism
\[
\tag{2}
\left\{\begin{gathered}
\diagram{\begin{tikzcd}[ampersand replacement=\&,column sep=large,row sep=small]
B'_0\times T_\rho\times T_\sigma\arrow[r]\&\mathcal G_{\rho,\sigma}\\
(\underline u,\underline a,\underline b)\arrow[r,mapsto]\&(\underline u\underline a^{-1},\underline u\underline b^{-1},\unsure{\nu|\delta}(\underline u))
\end{tikzcd}}
\end{gathered}\right.
\]
(where the fact that $\nu(\underline u)=\det(\underline u)-1$ is not of capital importance for the moment).

\margin{We now start from $\rho,\sigma,B'_0,\varepsilon_1$ (hence $L_1$) and we ask the question under what conditions, one can have a parametrization (2).}

\deleted{For} \inserted{To say} that this morphism \unsure{from the first|from the first ones} \unsure{into} $G\times G\times E'$ \deleted{gives} \inserted{maps} indeed into the solutions of (1), \deleted{it is necessary and sufficient} \inserted{\unsure{means}} also that we have identically
\[
\underline u\rho(\underline u)^{-1}=\underline u\sigma(u)^{-1}\varepsilon_1(\underline\nu(u))
\]
which can also be written, canceling $\unsure{\underline u|=}$ and taking inverses
\[
\tag{3}
\rho(\underline u)=\varepsilon_1(-\underline\nu(u))\sigma(\underline u).
\]
The possibility of finding $\underline\nu:B'_0\longrightarrow \mathbb G_a$ (morphism of schemes) \unsure{will give|would give}, via (2), a morphism $B'_0\times T_\rho\times T_\sigma\longrightarrow\mathcal G_{\rho,\sigma}\overset{\mathrm{dfn}}{=}$ scheme

\mspage{367}

% Printed archive footer outside the manuscript frame:
% 367
% https://grothendieck.umontpellier.fr
% Université de Montpellier, UMR CNRS 5149, Case courrier 051, Place Eugène Bataillon
% 34095 Montpellier cedex 5 - France

\scanpage{369}

of solutions of (1), is therefore equivalent to that
\deleted{\unsure{that}} \deleted{for every} $\sigma$ and $\rho$
\deleted{\unsure{have}} the same action
on $B'_0$, \uline{modulo} multiplication on the \unsure{left} by
\deleted{$\unsure{L_1}$} the subgroup \deleted{$\unsure{L_1}$}
\inserted{\unsure{unipoten\textsuperscript{t}}} with one parameter $L_1$ which
enters into the very writing of the equation (1) we are studying.
This is hard to conceive unless one admits that
$L_1\subset\rho(B'_0)\overset{\mathrm{df}}=B'_1$, which moreover implies
$\sigma(B'_0)\subset B'_1$ (and by a known argument $\sigma(B'_0)=B'_1$)
and $L_1=(B'_1)_u$. Conversely, if we have
$\rho(B'_0)=\sigma(B'_0)$, \deleted{\unsure{then for every}}
\inserted{i.e. $\sigma^{-1}$}\deleted{$\unsure{\rho(-)}$}$\rho$ normalizes $B'_0$,
\deleted{then} i.e. $\sigma^{-1}\rho$ a section of $B_0$, then
$\sigma(u)=\sigma u\sigma^{-1}$ and $\rho(u)=\rho u\rho^{-1}$ are two sections
of $B'_1$ which have the same determinant -- \unsure{hence abelianizatio\textsuperscript{n}},
hence which do coincide modulo $(B'_1)_u$.

Thus we see that the essential property for obtaining solutions of (1) via (2),
is to have
\[
(4)\qquad (\sigma^{-1}\rho)(B'_0)=B'_0\qquad\text{i.e.}\quad
\sigma^{-1}\rho\in\Gamma B_0\, ..
\]
\margin{\uline{N.B.} This condition \unsure{(if} $\sigma^{-1}\rho$ regular,
\unsure{i.e.} $\rho$ and $\sigma$ \unsure{are not homothetic}
\unsure{in any fiber)}, practically determines $B_0$
-- \deleted{\unsure{determines}} \unsure{over a field, there is an}
\uline{\unsure{ambiguity}} \unsure{possible} \unsure{between|of order} \uline{two}
(the \uline{\unsure{two}} Borels containing \unsure{the element in question}).}

For uniqueness of the representation of a triple $(\alpha,\beta,\nu)$ in the form
$(\underline u\underline a^{-1},\underline u\underline b^{-1},\underline\nu(\underline u))$,
it is necessary \deleted{to assume} \inserted{and it} suffices that we have
\[
(5)\qquad T_\sigma\cap T_\rho\cap B'_0=\{1\},
\]
\margin{It suffices for example that we have $T_\rho\cap T_\sigma=\unsure{\mathbb \mathbb G_m}$
and $\mathbb \mathbb G_m\cap B'_0=\{1\}$.}
\deleted{which is also equivalent to \unsure{our having}
\[
(6)\qquad T_\sigma\cap\unsure{\rho} B'_0\unsure{=}
\]}
the indeterminacy is

\mspage{368}

% Printed archive footer outside the manuscript frame:
% 368
% https://grothendieck.umontpellier.fr
% Université de Montpellier, UMR CNRS 5149, Case courrier 051, Place Eugène Bataillon
% 34095 Montpellier cedex 5 - France

\scanpage{370}
\mspage{486}

in the form
\diagram{\begin{tikzcd}[ampersand replacement=\&]
(\underline u,\underline a,\underline b) \arrow[r,mapsto] \&
(\underline u l,\underline a l,\underline b l)
\end{tikzcd}},
with $l$ a section of $T_\rho\cap T_\sigma\cap B'_0$.
There remains the question of the surjectivity of (2) -- can every solution of (1)
indeed be written in the form
$(u a^{-1},\underline u b^{-1},\nu(\underline u))$ ?
It amounts to the same (if $\rho,\sigma$ are regular)
to say that \deleted{\unsure{it}} we have surjectivity in
\[
(6)\qquad
\diagram{\begin{tikzcd}[ampersand replacement=\&]
B'_0 \arrow[r] \& \mathcal X_{\rho,\sigma} \\
\underline u \arrow[r,mapsto] \&
([\underline u,\rho],[\underline u,\sigma],\underline\nu(\underline u))
\end{tikzcd}}
\]
where $\mathcal X_{\rho,\sigma}$ is the scheme of simultaneous solutions
\deleted{in} in \deleted{$\unsure{\underline\nu,}$}
$\zeta,\eta,\nu$ of
\[
(7)\qquad
\det\zeta=\det\eta=1,\quad
\operatorname{Tr}\zeta\rho=\operatorname{Tr}\rho,\quad
\operatorname{Tr}\eta\sigma=\operatorname{Tr}\sigma,\quad
\zeta=\eta\varepsilon_1(\nu).
\]

It was essentially to obtain this surjectivity that we had been led
in section (II) to make the draconian hypotheses
(pages\ 435',436) which are expressed, at the cost of an étale,
by the conditions \deleted{\unsure{that|for}}
\[
(8)\qquad
\left\{\begin{array}{ll}
\mathrm{a)} & \sigma^{-1}\rho\text{ a section of }(B_0)_u=L_0
\quad(\text{and not only of }B_0)\\
\mathrm{b)} & B_0\text{ and }B_1=\rho(B_0)=
\deleted{\unsure{\rho}}\inserted{\sigma}(B_0)
\text{ in ``general position'',}\\
&\text{i.e. }L_0\text{ and }L_1=\rho(L_0)=\sigma(L_0)
\text{ independent;}\\
&\text{(in each fiber)}
\end{array}\right.
\]
but it is not clear whether these conditions
(suggested by earlier calculations, and indeed giving rise to
\unsure{some|some} very pretty calculations) are really indispensable for the result of parametrizati-

\mspage{369}

% Printed archive footer outside the manuscript frame:
% 369
% https://grothendieck.umontpellier.fr
% Université de Montpellier, UMR CNRS 5149, Case courrier 051, Place Eugène Bataillon
% 34095 Montpellier cedex 5 - France

\scanpage{371}

In our context it is \inserted{\unsure{quite}} natural to assume
$\varepsilon_0=\unsure{\sigma^{-1}\rho|-\sigma^{-1}\rho}$ unipotent
(or at least \unsure{quasi-unipotent}, i.e. \deleted{\unsure{a par}}
an iterate $\varepsilon_0^n$ unipotent) but condition b) of (8), more subtle,
appears less imperative. The sequel will show whether this restriction is judicious.

In the case where one considers an equation
\[
(9)\qquad
\underbrace{[\alpha,\rho]}_{\zeta}=
\varepsilon\underbrace{[\beta,\sigma]}_{\unsure{\eta}}
\varepsilon_1(\nu)
\qquad\deleted{\unsure{\varepsilon}}
\]
with a parameter $\varepsilon$ (which can in particular take the important value $-1$)
$\in\Gamma\mu_2$, \unsure{one|where} we have for the first time in these calculations
destroyed the symmetry between $\rho$ and $\sigma$, by introducing the condition
$\operatorname{Tr}\sigma=0$ as a condition of stability of the scheme
\inserted{$\mathcal X_\sigma$} of the $\eta$ satisfying
\[
\det\eta=1,\qquad\operatorname{Tr}\eta\sigma=\operatorname{Tr}\sigma
\]
under $\mu_2$. The essential result we were after, was the isomorphism
\[
(10)\qquad
\diagram{\begin{tikzcd}[ampersand replacement=\&,column sep=large]
B'_0\times T_\rho\times N'_\sigma \arrow[r,"\sim"] \& \mathcal G_{\rho,\sigma}^{*}\\
(\underline u,\underline a,\underline b) \arrow[r,mapsto] \&
(\alpha=\underline u\underline a^{-1},\beta=\underline u\underline b^{-1},
\mu=\det\underline u,\varepsilon=\varepsilon(\underline b))
\end{tikzcd}}
\]
where $N'_\sigma\subset N_\sigma$ is now the normalizer of $T_\sigma$
-- defined by the condition
\[
(11).\qquad \underline b\,\sigma(\underline b)^{-1}=
\varepsilon(\underline b)\deleted{\unsure{\sigma}}
\qquad\text{i.e.}\quad\underline b(\sigma)=\varepsilon(\underline b)\sigma
\]

\mspage{370}

% Printed archive footer outside the manuscript frame:
% 370
% https://grothendieck.umontpellier.fr
% Université de Montpellier, UMR CNRS 5149, Case courrier 051, Place Eugène Bataillon
% 34095 Montpellier cedex 5 - France

\scanpage{372}
\mspage{487}

That the morphism (10) does map into $\mathcal G_{\rho,\sigma}^{*}$,
is trivial by definition of $N'_\sigma$ and via the analogous fact for (2)
-- without a hypothesis of the kind $\operatorname{Tr}\sigma=0$.
But the possibility of obtaining (\unsure{local\textsuperscript{ly}})
any $\varepsilon$ via a $\underline b\in\Gamma N'_\sigma$
-- or simply $\varepsilon=-1$, if 2 invertible) requires precisely
$\operatorname{Tr}\sigma=0$.

Of course, we are going to consider $\mathcal G_{\rho,\sigma}^{*}$ as a group scheme,
by transport of \unsure{structure} from the product group structure of
$B'_0\times T_\rho\times N'_\sigma$.
We single out in it the subgroup
\[
(12)\qquad \mathcal G_{\rho,\sigma}^{\unsure{\flat|b}*}(0)
\quad\text{corresponding to :}\quad T'_0\times T_\rho\times N'_\sigma
\]
where
\[
(13)\qquad T_0=B'_0\cap B'_1
\]
is the distinguished maximal torus of $B'_0$, which
\deleted{for our} the canonical pinning is described as
\[
(14)\qquad T_0=\left\{\left(\begin{array}{cc}\mu&0\\0&1\end{array}\right.\unsure{\big\}|\big)}
\,\middle|\,\mu\in\Gamma\mathbb \mathbb G_m\right\}.
\]
We likewise set
\[
(15)\qquad \mathcal G_{\rho,\sigma}^{\unsure{\flat|b}}(0)
=\mathcal G_{\rho,\sigma}\cap\mathcal G_{\rho,\sigma}^{*\unsure{\flat|b}}(0),
\quad\text{corresponding to : }T'_0\times T_\rho\times T_\sigma
\]

\mspage{371}

% Printed archive footer outside the manuscript frame:
% 371
% https://grothendieck.umontpellier.fr
% Université de Montpellier, UMR CNRS 5149, Case courrier 051, Place Eugène Bataillon
% 34095 Montpellier cedex 5 - France